\documentclass[../../../include/open-logic-section]{subfiles}

\begin{document}

% Part: computability
% Chapter: computability-theory
% Section: applications-fixed-point

\olfileid{cmp}{thy}{apf}
\olsection{应用不动点定理}

不动点定理本质上允许我们利用其指标来定义部分可计算函数。
例如，我们可以找到一个指标~$e$，使得对每个~$y$，
\[
\cfind{e}(y) = e + y.
\]
再举一个例子，可以利用不动点定理的证明，来设计一个打印自身的 Java 或 C++ 程序。

记住，如果对每个 $e$，我们令 $W_e$ 为 $\cfind{e}$ 的定义域，那么序列 $W_0$, $W_1$, $W_2$,~\dots 就枚举了所有的可计算枚举集。
这些集合中有些是可计算的。
可以问，是否存在一个算法，它以值 $x$ 为输入，并在 $W_x$ 恰好可计算时，返回其特征函数的一个指标。
答案是否定的，不存在这样的算法：

\begin{thm}
不存在满足下述性质的部分可计算函数 $f$：每当 $W_e$ 可计算时，$f(e)$ 有定义且 $\cfind{f(e)}$ 是其特征函数。
\end{thm}

\begin{proof}
设 $f$ 为任意可计算函数；我们将构造一个 $e$，使得 $W_e$ 可计算，但 $\cfind{f(e)}$ 不是其特征函数。
利用不动点定理，我们可以找到一个指标 $e$，使得
\[
\cfind{e}(y) \simeq
\begin{cases}
0 & \text{若 $y=0$ 且 $\cfind{f(e)}(0) \fdefined = 0$} \\
\text{无定义} & \text{否则。}
\end{cases}
\]
也就是说，$e$ 是将不动点定理应用于下式定义的函数而得到的：
\[
g(x,y) \simeq
\begin{cases}
0 & \text{若 $y=0$ 且 $\cfind{f(x)}(0) \fdefined = 0$} \\
\text{无定义} & \text{否则。}
\end{cases}
\]
直观上看，我们可以得出 $g$ 是部分可计算的，如下所示：对输入 $x$ 和 $y$，算法首先检查 $y$ 是否等于~$0$。
若等于，算法计算 $f(x)$，然后用通用机计算 $\cfind{f(x)}(0)$。
如果该计算停机并返回~$0$，算法就返回~$0$；否则，算法不停机。

但现在请注意，如果 $\cfind{f(e)}(0)$ 有定义且等于 $0$，那么 $\cfind{e}(y)$ 有定义当且仅当 $y$ 等于 $0$，所以 $W_e = \{ 0 \}$。
如果 $\cfind{f(e)}(0)$ 无定义，或有定义但不等于 $0$，那么 $W_e = \emptyset$。
无论哪种情况，$\cfind{f(e)}$ 都不是~$W_e$ 的特征函数，因为它在输入 $0$ 上给出了错误的值。
\end{proof}

\end{document}